\documentclass{article}
\begin{document}
\title{Evidence-Grounded Study of SQLite index tradeoffs on the UCI Bike Sharing hourly workload}
\maketitle

\section{Abstract}
We study SQLite index tradeoffs on the UCI Bike Sharing hourly workload with a local, auditable experiment pipeline. The best kept trial was `ablation-season-index-seed3` with primary metric 32.412629149999994.

\section{Introduction}
The goal is to convert a research idea into a paper candidate without losing provenance.
This draft only reports evidence that exists in the run artifacts.

\section{Related Work}
The screened literature includes MOSIQS: Persistent Memory Object Storage With Metadata Indexing and Querying for Scientific Computing \cite{khan2021mosiqs}, AirIndex: Versatile Index Tuning Through Data and Storage \cite{chockchowwat2023airindex}, DART \cite{zhang2018dart}, mCerebrum \cite{hossain2017mcerebrum}, Benchmarking PHP–MySQL Communication: A Comparative Study of MySQLi and PDO Under Varying Query Complexity \cite{andrijevi2025benchmarking}, TASM: A Tile-Based Storage Manager for Video Analytics \cite{daum2021tasm}, Index of SQLite Commands \cite{paper2019index}, Database Workload Optimization \cite{fritchey2018database}, An adaptive path index for XML data using the query workload \cite{min2005adaptive}, AI-Based Automated SQL Query Generation for SQLite Databases in Mobile Forensics \cite{pawlaszczyk2026aibased}, snowquery: SQL Interface to 'Snowflake', 'Redshift', 'Postgres', 'SQLite', and 'DuckDB' \cite{mermelstein2023snowquery}, Erratum to: “An adaptive path index for XML data using the query workload” \cite{min2006erratum}, Database Workload Optimization \cite{paper2009database}, Database Workload Optimization \cite{fritchey2014database}, Text Document Annotation and Retrieval Based on Content of the Document and Query Workload \cite{paper2016text}, Introducing SQLite \cite{paperndintroducing}, Driving through the Network: Performance and Workload under Latency and Video Impairments \cite{trautmannsheimer2026driving}, ATTI: Workload-Aware Query Adaptive OcTree Based Trajectory Index \cite{meng2013atti}, Query Workload Driven Summarization for P2P Query Routing \cite{nguyen2008query}, Database Workload Optimization \cite{fritchey2012database}, SQLite Internals \cite{paperndsqlite}, Research on SQLite Database Query Optimization Based on Improved PSO Algorithm \cite{zhao2016research}, A Cost-Effective Query Optimizer for Multi-Tenant Parallel RDBMSs Leveraging Workload Prediction \cite{danaouindcosteffective}, SQLite Internals and New Features \cite{allen2010sqlite}, adbcsqlite: 'Arrow' Database Connectivity ('ADBC') 'SQLite' Driver \cite{dunnington2023adbcsqlite}, Analyzing SQLite Databases \cite{languedoc2016analyzing}, Decentralized, Energy-Efficient, Low Latency and Less Homogeneous Settings based Workload Management in Enterprise Clouds \cite{bhuvaneshwari2016decentralized}, Using SQLite with PHP \cite{feiler2015using}, Parallel selection query processing involving index in parallel database systems \cite{rahayundparallel}, Index selection \cite{sun2013index}, A Copula-Based Sample Selection Binary Choice Model for Difference Analysis Among Private Bike and Bike Sharing in Lyon (France) \cite{havet2024copulabased}, Index \cite{paper2015index}, Large-Scale Dockless Bike Sharing Repositioning Considering Future Usage and Workload Balance \cite{hua2022largescale}, Bike Sharing Systems \cite{paper2012bike}, Dynamic Workload-Aware Bike Rebalancing for Bike-Sharing Systems \cite{luo2023dynamic}, Analyzing Bike Repositioning Strategies Based on Simulations for Public Bike Sharing Systems: Simulating Bike Repositioning Strategies for Bike Sharing Systems \cite{wang2013analyzing}, Bike Sharing Systems \cite{brinkmann2020bike}, Bikeability Index in Bike-Sharing Systems: A Dual-Level Assessment Integrating Station Accessibility and Cycling Environment \cite{zhang2026bikeability}, The impact of the introduction of e-bike sharing on the usage of bike sharing \cite{li2023impact}, Investigation on the impact of new bike stations on a bike-share system based on a complex bike-sharing network \cite{kim2023investigation}, Shared Bike Demand Prediction by Using Metro and Bike Sharing Networks’ Features \cite{sadeghraimoghaddam2024shared}, Service Network Design of Bike Sharing Systems \cite{vogel2016service}, Bike Sharing Atlas: Visual Analysis of Bike-Sharing Networks \cite{oppermann2018bike}, Bike Sharing in the Context of Urban Mobility \cite{vogel2016bike}, Large-scale dockless bike sharing repositioning considering future usage and workload balance \cite{hua2022largescale}, Impacts of Bike Sharing on Transit Ridership \cite{aljerindimpacts}, Figure 4: The sliding window technique for predicting hourly bike sharing demand. \cite{paperndfigure}, Station-Level Hourly Bike Demand Prediction for Dynamic Repositioning in Bike Sharing Systems \cite{wu2019stationlevel}, Framework for Hourly Demand Forecasting of Bike-Sharing Stations: Case Study of the Four Main Gate Areas in Seoul \cite{hong2022framework}, Correction: Graph convolutional network approach applied to predict hourly bike-sharing demands considering spatial, temporal, and global effects \cite{kim2022correction}, Graph convolutional network approach applied to predict hourly bike-sharing demands considering spatial, temporal, and global effects \cite{kim2019graph}, Spatiotemporal Data-Driven Hourly Bike-Sharing Demand Prediction Using ApexBoost Regression \cite{biswas2025spatiotemporal}, Analyzing Tail Latency in Serverless Clouds with STeLLAR \cite{ustiugov2021analyzing}, Road-Specific Exploration of Bike-Sharing Usage Changes after Construction of Bike Lanes \cite{li2022roadspecific}, From System-Wide to Road-Specific Exploration of Bike Trips for Changes in Bike Sharing System Usage after Construction of Bike Lanes \cite{li2022systemwide}, Tradeoffs between power management and tail latency in warehouse-scale applications \cite{kanev2014tradeoffs}, Research Department - Prices \&amp; Statistics - Price Indexes - Wage Indexes - Minimum Hourly Rates of Wages by Industrial Groups - Correspondence, Memoranda and Blue Sheets - 1950 - 1959 \cite{paper2022research}, The bike sharing rebalancing problem: Mathematical formulations and benchmark instances \cite{dellamico2014bike}, Query similarity index based query preprocessing mechanism for multiapplication sharing wireless sensor networks \cite{verma2020query}, Figure 6: Proposed MLP to predict future bike sharing demand. \cite{paperndfigure}.

\section{Method}
The configured domain experiment workspace executes the prespecified trials.
Each trial writes structured metrics, and the pipeline keeps only metric improvements.

\section{Experiments}
\noindent `baseline-seed0`: metric=46.5152229, decision=keep, status=ok.\\
\noindent `baseline-seed1`: metric=38.69099614999999, decision=keep, status=ok.\\
\noindent `baseline-seed2`: metric=45.3883545, decision=discard, status=ok.\\
\noindent `baseline-seed3`: metric=53.22321424999999, decision=discard, status=ok.\\
\noindent `ablation-hour-index-seed1`: metric=42.04770625, decision=discard, status=ok.\\
\noindent `ablation-weather-index-seed2`: metric=49.67862115, decision=discard, status=ok.\\
\noindent `ablation-season-index-seed3`: metric=32.412629149999994, decision=keep, status=ok.\\
\noindent `ablation-composite-index-seed4`: metric=35.3873705, decision=discard, status=ok.\\

\section{Results}
The best kept trial was `ablation-season-index-seed3` with primary metric 32.412629149999994.
\title{Research Decision}
\maketitle

Decision: PROCEED

Proceed with trial `ablation-season-index-seed3` as the current evidence baseline.

\section{Limitations}
Claims are limited to the registered assets, evaluation units, trials, metrics, and compute budget recorded in this run.

\section{Conclusion}
The workflow now links literature, experiment metrics, and paper text through auditable artifacts.

\section{Revision Notes}
The revision preserves only screened citations and verified experiment metrics. Open reviewer concerns are tracked in `reviews.md`.

<!-- Review summary preserved for audit -->
<!-- \# Reviews

\section{Reviewer A: Methodology}
\noindent The paper is structurally complete and reports the current evidence.\\
\section{Reviewer B: Evidence}
\noindent Claims should remain tied to verified metrics and screened citations.\\
\section{Reviewer C: Venue Readiness}
\noindent Venue readiness remains conditional on exact template and policy checks.\\
 -->
\end{document}
